\chapter*{Problem Statement}
\noindent

The reliability of digital memory systems is under growing threat. As NAND flash cells are operated beyond their endurance limits in multi-level cell (MLC) and triple-level cell (TLC) configurations, raw bit error rates (BER) reach $10^{-3}$ to $10^{-2}$ per bit per read cycle near end-of-life. DRAM operating at reduced supply voltage for energy efficiency exhibits BERs in the range of $10^{-6}$ to $10^{-4}$ per cell per hour. In a representative 4,096-bit data block at 5\% BER, approximately 200 bits are expected to be flipped---far beyond the correction capacity of typical on-chip ECC engines, which are sized for 1--8 errors.

The dominant deployed ECC technology, BCH codes, addresses bit-flip correction algebraically. A BCH code correcting $t$ errors in an $n$-bit block requires approximately $2t \cdot \lceil \log_2(n+1) \rceil$ parity bits and a decoder implementing the Berlekamp-Massey algorithm over the Galois field $\text{GF}(2^m)$. For $t = 200$ errors and $n = 4{,}096$ bits, this yields approximately 4,800 parity bits (117\% overhead in the best case), rising to $\sim$174\% when $t$ is conservatively sized for the 95th percentile of the error distribution---a necessary practice because BCH provides \emph{no} correction beyond its rated $t$. This ``honest-$t$'' penalty doubles or triples overhead in practice.

Beyond overhead, BCH codes impose two structural limitations that are increasingly mismatched to modern deployment contexts. First, the Berlekamp-Massey decoder requires Galois field multiplications, which demand either dedicated silicon area for finite-field arithmetic units or slow software implementations---neither of which is acceptable in area-constrained memory controllers or latency-sensitive inference engines. Second, BCH codes assume an \emph{independent, random} error model. Contiguous burst errors---caused by flash cell-to-cell interference, DRAM row-hammer disturbance, or memory controller faults---concentrate many errors in a small region, requiring $t$ proportional to the burst length $b$ to correct. Without interleaving (which adds latency and buffer overhead), a burst of 50 contiguous flips requires a $t \geq 50$ BCH code for the entire block, multiplying overhead for all data even when burst errors are rare.

The rise of deep neural network (DNN) inference at the edge compounds these challenges. AI accelerators store large weight matrices in on-chip SRAM or off-chip DRAM and access them at high bandwidth and potentially reduced supply voltage. A single bit flip in the exponent of a 16-bit floating-point weight can multiply its value by up to 65,536, causing catastrophic accuracy collapse in the affected layer. At the same time, the inherent noise tolerance of neural networks means that a high success rate of \emph{approximate} correction---not algebraic certainty---is sufficient for reliable inference. No deployed ECC scheme is designed to exploit this tolerance.

\noindent\underline{The absence of a practical, hardware-friendly error correction scheme that simultaneously achieves low metadata overhead, native burst resilience, and applicability to high bit error rates without Galois field arithmetic represents a critical gap that limits the reliability of modern memory-intensive and AI computing systems.}
